\chapter{Bounds on thermal traces, and seeds from trace tables}
\label{chap:traces}

The finite inputs in \cite{OpenAIHaldane, OpenAIBoundary} come from short-chain traces and
variational bounds. This chapter records the generic tools that turn exact integer computations
into those inputs. The spin-1 instances follow in Chapters~\ref{chap:heisenberg}
and~\ref{chap:open}; the kernel computations appear in Chapter~\ref{chap:seedkernel}.

\section{Spectral bounds}

Quadratic forms bound spectra without a numerical eigensolve.
Counting bounds also give lower and upper bounds on thermal traces.

\begin{definition}[Quadratic forms and spectral counting]
  \label{def:spectral-forms}
  \lean{CertifiedGaps.Spectral.qform, CertifiedGaps.Spectral.nsq,
    CertifiedGaps.Spectral.eigCount}
  \leanok
  For a finite Hermitian matrix \(A\), set
  \(q_A(x)=\operatorname{Re}\langle x,Ax\rangle\), \(N(x)=\langle x,x\rangle\), and
  \(N_A(t)=\#\{i:\lambda_i(A)<t\}\), counting multiplicities.
\end{definition}

\begin{lemma}[Operator and functional-calculus bounds]
  \label{lem:spectral-calculus}
  \lean{CertifiedGaps.Spectral.abs_qform_le_of_schur,
    CertifiedGaps.Spectral.abs_eigenvalues_le_one_of_qform,
    CertifiedGaps.Spectral.nsq_hermFun_mulVec_le,
    CertifiedGaps.Spectral.eigenvalues_mem_of_prod_eq_zero}
  \leanok
  \uses{def:spectral-forms}
  Suppose \(|A_{ij}|\le K_{ij}=K_{ji}\), \(q_i>0\), and
  \(\sum_j K_{ij}q_j\le c q_i\). Then \(|q_A(x)|\le cN(x)\).
  For Hermitian \(A\), \(|q_A(x)|\le N(x)\) for all \(x\) implies
  \(|\lambda_i(A)|\le1\). If \(|f(\lambda_i(A))|\le\eta\) for every \(i\), then
  \(N(f(A)x)\le\eta^2N(x)\).
  If \(\prod_{r\in R}(A-rI)=0\) for a finite list \(R\) of real numbers, each eigenvalue lies in
  that list.
\end{lemma}
\begin{proof}
  \leanok
  Apply the weighted Schur bound and diagonalize the Hermitian matrix.
\end{proof}

\begin{lemma}[Counting by subspaces and exact inertia]
  \label{lem:spectral-counting}
  \lean{CertifiedGaps.Spectral.eigCount_le_of_forall_le_qform,
    CertifiedGaps.Spectral.le_eigCount_of_forall_qform_lt,
    CertifiedGaps.Spectral.LDLCert, CertifiedGaps.Spectral.LDLCert.exists_neg_subspace}
  \leanok
  \uses{def:spectral-forms}
  If a subspace \(V\subseteq\C^d\) has codimension at most \(K\) and
  \(tN(x)\le q_A(x)\) on \(V\), then \(N_A(t)\le K\).
  If \(q_A(x)<tN(x)\) for every nonzero \(x\in V\), then \(\dim V\le N_A(t)\).
  A rational factorization \(T=LDL^{\mathsf T}\), with \(L\) unit lower triangular and \(D\)
  diagonal, supplies a negative subspace of dimension equal to the number of negative entries
  of \(D\).
\end{lemma}
\begin{proof}
  \leanok
  Intersect \(V\) with the spectral subspace on the other side of \(t\).
  For the inertia certificate, pull back the negative coordinate subspace through \(L^{\mathsf T}\).
\end{proof}

\begin{lemma}[Schur correction for a trial subspace]
  \label{lem:spectral-schur}
  \lean{CertifiedGaps.Spectral.eigCount_le_of_schur_LDL}
  \leanok
  \uses{lem:spectral-counting}
  Let \(H\) be Hermitian, \(m>0\), and
  \(q_H(u)\ge(t+m)N(u)\) whenever \(W^*u=0\).
  Suppose
  \(\|(H-tI)Z-WB\|_{\mathrm F}^2\le m\delta\) and
  \(\operatorname{range}Z+\ker W^*=\C^d\).
  Let a rational matrix \(T=LDL^{\mathsf T}\) satisfy
  \(q_T(v)\le q_{Z^*(H-tI)Z-\delta I}(v)\) for every \(v\).
  Then \(N_H(t)\) is at most the number of negative entries of \(D\).
  This upper bound does not require \(L\) to be invertible.
\end{lemma}
\begin{proof}
  \leanok
  Bound the cross term by the residual norm and \(m\), then use the nonnegative subspace
  supplied by the factorization.
\end{proof}

\begin{lemma}[Thermal sums from counting bounds]
  \label{lem:spectral-trace-bins}
  \lean{CertifiedGaps.Spectral.sum_le_binned, CertifiedGaps.Spectral.binned_le_sum,
    CertifiedGaps.Spectral.antitone_exp_neg_mul}
  \leanok
  \uses{def:spectral-forms, lem:spectral-counting}
  Let \(f\) be nonincreasing, \(t_k\) nondecreasing, and \(e_i\ge t_0\) for \(i\) in a finite
  set \(S\).
  If \(\#\{i\in S:e_i<t_{k+1}\}\le N_{k+1}\) for \(k<m\), then
  \[
    \sum_{i\in S}f(e_i)\le
    \sum_{k<m}N_{k+1}\bigl(f(t_k)-f(t_{k+1})\bigr)+|S|f(t_m).
  \]
  If \(f\ge0\) and \(M_k\le\#\{i\in S:e_i<t_k\}\) for \(k<m\), then
  \(\sum_{k<m}M_k(f(t_k)-f(t_{k+1}))\le\sum_{i\in S}f(e_i)\), without the assumption \(e_i\ge t_0\).
  In particular, \(f(x)=e^{-\beta x}\) is admissible for \(\beta\ge0\).
\end{lemma}
\begin{proof}
  \leanok
  Telescope the increments of \(f\) and count the levels below each threshold.
\end{proof}

\section{Poisson truncation and Horner recursions}

The thermal calculations in \cite[Appendices A and B]{OpenAIHaldane} approximate the square root
of a shifted Gibbs operator. Integer division introduces a bounded error at each Horner step.

\begin{lemma}[Normalized Poisson polynomial]
  \label{lem:trace-poisson}
  \lean{CertifiedGaps.Trace.poissonWeight, CertifiedGaps.Trace.poissonError,
    CertifiedGaps.Trace.abs_exp_sub_poissonWeight_le,
    CertifiedGaps.Trace.nsq_exp_sub_poisson_mulVec_le,
    CertifiedGaps.Trace.poissonError_98_280, CertifiedGaps.Trace.poissonError_84_256}
  \leanok
  \uses{lem:spectral-calculus}
  For rational \(m\ge0\), put
  \[
    p_j=\frac{m^j/j!}{\sum_{\ell=0}^J m^\ell/\ell!},
    \qquad \eta=2e^{-m}\sum_{j>J}\frac{m^j}{j!}.
  \]
  For \(|x|\le1\), \(|e^{m(x-1)}-\sum_{j=0}^Jp_jx^j|\le\eta\).
  The same error in operator norm holds for a finite Hermitian contraction.
  The checked constants are \(\eta(98,280)<2^{-84}\) and \(\eta(84,256)<2^{-116}\).
\end{lemma}
\begin{proof}
  \leanok
  Bound the omitted exponential series and the normalization error separately.
  Apply functional calculus to the scalar inequality.
\end{proof}

\begin{lemma}[Floored integer Horner columns]
  \label{lem:trace-horner}
  \lean{CertifiedGaps.Trace.hornerFloor, CertifiedGaps.Trace.intDiv_eq_floor,
    CertifiedGaps.Trace.hornerFloor_error_card}
  \leanok
  Let \(D\) be an integer \(d\times d\) matrix, \(s>0\) an integer, and \(D/s\) a Euclidean
  contraction. Starting from \(X_{J+1}=0\), set
  \(X_j=\lfloor DX_{j+1}/s\rfloor+c_je_r\), with the floor taken entrywise.
  If \(|c_j-a_j|\le1\) for \(0\le j\le J\), then
  \[
    \left\|X_0-\sum_{j=0}^J a_j(D/s)^je_r\right\|_2
    \le(J+1)(\sqrt d+1).
  \]
  Integer division is floor division, including for negative entries.
\end{lemma}
\begin{proof}
  \leanok
  Each step contributes at most \(\sqrt d\) from division and \(1\) from the coefficient.
  The contraction does not increase earlier errors.
\end{proof}

\begin{lemma}[Horner columns over Eisenstein integers]
  \label{lem:trace-eisenstein}
  \lean{CertifiedGaps.Trace.EisensteinInt, CertifiedGaps.Trace.hornerFloorE,
    CertifiedGaps.Trace.hornerFloorE_error_nsq}
  \leanok
  \uses{def:spectral-forms}
  Write \(\omega=e^{2\pi i/3}\) and use pairs of integers for \(a+b\omega\).
  Let \(D/s\) be a complex Euclidean contraction with entries in \(\mathbb Z[\omega]/s\),
  where \(s>0\). Floor the two integer coordinates at each step.
  With real coefficients \(a_j\), integer \(c_j\), and \(|c_j-a_j|\le1\), the squared column
  error is at most \(((J+1)(2\delta+1))^2\) for any \(\delta\ge\sqrt d\).
\end{lemma}
\begin{proof}
  \leanok
  The two coordinate errors have complex norm at most \(2\).
  Propagate this bound through the contraction.
\end{proof}

\begin{lemma}[Weighted Horner error]
  \label{lem:trace-weighted}
  \lean{CertifiedGaps.Trace.IsWeightedContraction,
    CertifiedGaps.Trace.hornerFloor_error_weighted,
    CertifiedGaps.Trace.isWeightedContraction_of_similar,
    CertifiedGaps.Trace.gibbsTrace_enclosure_weighted}
  \leanok
  \uses{lem:trace-horner, lem:trace-poisson}
  For weights \(w_i\ge0\), suppose \(D/s\) contracts the norm
  \(\|x\|_w^2=\sum_i w_i x_i^2\), with \(s>0\).
  Under the coefficient assumptions of Lemma~\ref{lem:trace-horner}, the error in column \(r\)
  is at most \((J+1)(\delta+\gamma)\), where
  \(\delta\ge\sqrt{\sum_iw_i}\) and \(\gamma\ge\sqrt{w_r}\).
  If \(D=s\,\operatorname{diag}(j)A\operatorname{diag}(j)^{-1}\), \(j_i>0\), and \(A\) is a
  Hermitian contraction, take \(w_i=j_i^{-2}\).
  Undo this similarity before summing squared columns. Suppose
  \(e^{-b(H+cI)/2}=e^{m(A-I)}\), \(Q>0\), \(|c_k-Qp_k|\le1\) for the Poisson weights \(p_k\),
  \(j_i\le j_{\max}\), \(\gamma\ge j_i^{-1}\), and
  \(\epsilon\ge\sqrt d\,(j_{\max}(J+1)(\delta+\gamma)+Q\eta)\). With
  \(W=\sum_{r,w}(j_r/j_w)^2X_r(w)^2\),
  \[
    \frac{e^{bc}\max(0,\sqrt W-\epsilon)^2}{Q^2}\le\Tr e^{-bH}
    \le\frac{e^{bc}(\sqrt W+\epsilon)^2}{Q^2}.
  \]
\end{lemma}
\begin{proof}
  \leanok
  Bound each rounding error in the weighted norm, then apply the Poisson error and the
  Frobenius norm triangle inequality.
\end{proof}

\begin{lemma}[Packed evaluation preserves the integer recurrence]
  \label{lem:trace-packed}
  \lean{CertifiedGaps.Trace.Packed.pack_lane, CertifiedGaps.Trace.Packed.Band.check_spec,
    CertifiedGaps.Trace.Packed.hornerFloor_eq_specCol,
    CertifiedGaps.Trace.Packed.run_spec, CertifiedGaps.Trace.Packed.run_numerator,
    CertifiedGaps.Trace.Packed.runCols_spec, CertifiedGaps.Trace.Packed.runCols_numerator}
  \leanok
  \uses{lem:trace-horner}
  Suppose the packing parameters satisfy \texttt{Params.Ok}, a band matrix has its indices
  in range and satisfies \texttt{RowsOk}, and every coefficient is at most the offset \(B\).
  A successful packed run decodes to the integer Horner columns it was asked for: a window of
  consecutive columns, or any list of distinct columns.
  For row and column weight lists of the specified lengths, its numerator equals the weighted
  sum of their squared entries.
  The Boolean band check implies the two band conditions.
\end{lemma}
\begin{proof}
  \leanok
  Extract each radix lane, verify the multiply-and-shift division, and induct over the
  coefficient list. The range and offset bounds prevent carries from changing another lane.
\end{proof}

\section{Orbit representatives and twisted traces}

Column symmetries reduce a trace computation to representatives with exact multiplicities.
The symmetry can carry a phase or complex conjugation.

\begin{lemma}[Column norms on symmetry orbits]
  \label{lem:trace-orbits}
  \lean{CertifiedGaps.Trace.IsColumnSymmetry, CertifiedGaps.Trace.IsConjColumnSymmetry,
    CertifiedGaps.Trace.nsq_gibbsHalf_rep,
    CertifiedGaps.Trace.nsq_gibbsHalf_rep_of_norms,
    CertifiedGaps.Trace.IsConjColumnSymmetry.nsq_gibbsHalf}
  \leanok
  Let \(F=e^{-b(H+cI)/2}\).
  A unitary symmetry of \(H\) that sends \(e_r\) to a unit phase times \(e_{\sigma(r)}\)
  preserves \(\|Fe_r\|^2\). For Hermitian \(H\), the same holds for an antiunitary column symmetry.
  If each generator preserves these norms and \(\mathrm{rep}(r)\) belongs to the generated
  orbit of \(r\), then \(\|Fe_r\|^2=\|Fe_{\mathrm{rep}(r)}\|^2\).
\end{lemma}
\begin{proof}
  \leanok
  Commute the symmetry through the exponential. Norm preservation is closed under composition.
\end{proof}

\begin{lemma}[Rational trace bounds from integer column totals]
  \label{lem:trace-frobenius}
  \lean{CertifiedGaps.Trace.gibbsTrace_eq_exp_mul_sum_nsq,
    CertifiedGaps.Trace.gibbsTrace_enclosure_of_intColumns,
    CertifiedGaps.Trace.gibbsTrace_enclosure_of_rep}
  \leanok
  \uses{def:spectral-forms}
  Let \(H\) be Hermitian and \(F=e^{-b(H+cI)/2}\).
  Suppose representative columns have multiplicities \(\nu_\rho\) that reproduce
  \(\sum_r\|Fe_r\|^2\), and integer columns \(X_\rho\) satisfy
  \(\|X_\rho-QFe_\rho\|\le E\), where \(Q>0\), \(E\ge0\).
  Set \(W=\sum_\rho\nu_\rho\|X_\rho\|^2\), \(r=\lfloor\sqrt W\rfloor\), and suppose
  \(\mathcal E\ge0\), \(\sum_\rho\nu_\rho E^2\le\mathcal E^2\), and
  \(s\le e^{bc}\le s+t\). Then
  \[
    \frac{s\max(0,r-\mathcal E)^2}{Q^2}
    \le \Tr e^{-bH}\le
    \frac{(s+t)(r+1+\mathcal E)^2}{Q^2}.
  \]
  This interval is an enclosure of the trace; an enclosure of a number is an interval proved
  to contain it.
\end{lemma}
\begin{proof}
  \leanok
  Use \(\Tr e^{-bH}=e^{bc}\sum_r\|Fe_r\|^2\), the weighted Frobenius triangle inequality,
  and \(r^2\le W<(r+1)^2\).
\end{proof}

\section{Seeds and reseeding}

Twisted traces separate symmetry sectors. Bounds on their moments then control the leading
eigenvalue and the remaining mass, as in \cite[Section 5]{OpenAIHaldane}.

\begin{definition}[Sector schemes and finite approximations]
  \label{def:seed-scheme}
  \lean{CertifiedGaps.SectorScheme, CertifiedGaps.SectorScheme.Valid,
    CertifiedGaps.SectorScheme.d2, CertifiedGaps.SectorScheme.ion,
    CertifiedGaps.SectorApprox}
  \leanok
  A sector scheme records integer coefficients \(C_{gs}\) and a rational left inverse.
  Validity requires that the rational matrix be a left inverse of \(C\), \(C_{00}=1\), and
  \(C_{0s}\ge0\).
  Its traces are \(Z_g(n)=\sum_s C_{gs}\sum_{i\in s}x_i^n\).
  A sector approximation means that, for every \(\delta>0\), such finite real power sums
  approximate all requested traces and lengths within \(\delta\).
  The \(D_2\) scheme is \((Z,Z_P)=(J+3N,J-N)\).
  The three-sector scheme is
  \((Z,Z_P,Z_C)=(I+2O+3N,I+2O-N,I-O)\).
\end{definition}

\begin{lemma}[Inverting a trace table]
  \label{lem:seed-table-inversion}
  \lean{CertifiedGaps.TraceBounds, CertifiedGaps.TraceBounds.Holds,
    CertifiedGaps.TraceBounds.HoldsStrict, CertifiedGaps.TraceBounds.momentBounds,
    CertifiedGaps.TraceBounds.holdsStrict_centered}
  \leanok
  \uses{def:seed-scheme}
  A trace table gives optional rational lower and upper bounds at specified lengths.
  For a valid sector scheme, applying its rational left inverse with the correct interval
  endpoints bounds each sector moment.
  If traces differ from their centers by at most \(\epsilon_t<\epsilon\), they satisfy the
  table with radius \(\epsilon\) strictly.
\end{lemma}
\begin{proof}
  \leanok
  In each linear combination use upper endpoints for positive coefficients and lower endpoints
  for negative coefficients.
\end{proof}

\begin{lemma}[Capped masses and a dominant sector]
  \label{lem:seed-capped}
  \lean{CertifiedGaps.sum_pow_le_capped, CertifiedGaps.SectorSeed,
    CertifiedGaps.SectorSeed.Valid, CertifiedGaps.SectorSeed.seed}
  \leanok
  If \(0\le t_i\le h\), \(\sum_i t_i\le m\), and \(h\le m\le2h\), then
  \(\sum_i t_i^r\le h^r+(m-h)^r\) for every integer \(r\ge1\).
  For a valid \texttt{SectorSeed}, suppose every specified eigenvalue cap holds strictly,
  every sector's \(q\)-th moment is at most its recorded mass, and
  \(1<Z_{2qa}\), where \(Z_n=\sum_s\mathrm{mult}_s\sum_{i\in s}x_i^n\).
  Then
  \((1-\es)Z_{qa}^2\le Z_{2qa}\le B_\tau\), with the defect and thermal bound recorded
  by the certificate.
\end{lemma}
\begin{proof}
  \leanok
  Concentrate a bounded mass at its cap. The validity inequalities force a dominant value
  in sector \(0\), then bound the remaining mass.
\end{proof}

\begin{theorem}[Purities from a valid trace-table certificate]
  \label{thm:seed-table}
  \lean{CertifiedGaps.TableCertificate, CertifiedGaps.TableCertificate.Valid,
    CertifiedGaps.TableCertificate.purities, CertifiedGaps.seed_limit,
    CertifiedGaps.purities_of_shifted, CertifiedGaps.RPRing.seeds_of_certificate}
  \leanok
  \uses{lem:seed-table-inversion, lem:seed-filter, lem:seed-capped}
  Let \(C\) be a valid table certificate, with \(n_0=qa\), inverse temperature \(\beta\),
  shift \(\kappa\), and defects \(\es,\et\). Assume the sector approximation at its requested
  lengths, the strict table bounds for
  \(\widetilde Z_g(n,\beta)=e^{-\beta\kappa n}Z_g(n,\beta)\), and \(Z_0(n_0,\beta)>0\).
  If \(\widetilde Z_0(2n_0,\beta)>1\) and
  \(\widetilde Z_0(2n_0,2\beta)>1\), then
  \(\Ps(n_0,\beta)\ge1-\es\) and \(\Pt(2n_0,\beta)\ge1-\et\).
\end{theorem}
\begin{proof}
  \leanok
  Apply the filter and capped-mass bounds to sufficiently close finite approximations.
  Pass the non-strict inequalities to the limit and cancel the energy shift.
\end{proof}

\begin{lemma}[Changing inverse temperature]
  \label{lem:seed-temperature-lift}
  \lean{CertifiedGaps.BoltzmannRep, CertifiedGaps.BoltzmannRep.lift,
    CertifiedGaps.boltzmann_sq_le_of_dominant}
  \leanok
  For a finite Boltzmann sum \(f(\beta)=\sum_i e^{-\beta E_i}\) and integers \(1\le j\le l\),
  \(f((l/j)\beta)^j\le f(\beta)^l\).
  More precisely, suppose one weight at \(\beta_0\) exceeds \(1\),
  \(f(\beta_0)\le q\), \(R\ge0\), and \((q-1)^l<R^j\). Then
  \[
    f((l/j)\beta_0)^2\le(1+R)^2 f(2(l/j)\beta_0).
  \]
\end{lemma}
\begin{proof}
  \leanok
  Compare powers of nonnegative weights, separating the dominant weight in the second bound.
\end{proof}

\begin{theorem}[A lifted three-sector seed]
  \label{thm:seed-lifted}
  \lean{CertifiedGaps.LiftedSeed, CertifiedGaps.LiftedSeed.Valid,
    CertifiedGaps.LiftedSeed.purities}
  \leanok
  \uses{thm:seed-table, lem:seed-temperature-lift}
  Let \(C\) be a valid lifted seed, \(k=qr\), and \(\beta'=(l/j)\beta\).
  Assume the three-sector approximations at \(\beta,\beta'\) on the certificate's two length
  sets and the strict shifted table at \(\beta\).
  Suppose each twisted trace at length \(k\) is a finite Boltzmann sum.
  Suppose also that the untwisted trace at length \(4k\) has a finite energy representation
  with some energy below \(-4k\kappa\), and \(Z_0(2k,\beta')>0\).
  Then \(\Ps(2k,\beta')\ge1-\es\) and \(\Pt(4k,\beta')\ge1-\et\).
\end{theorem}
\begin{proof}
  \leanok
  Lift the short-trace bounds in temperature and run the certificate's rational dominance
  steps in the three sectors.
\end{proof}

\begin{theorem}[Reseeding at another length]
  \label{thm:seed-reseed}
  \lean{CertifiedGaps.ReseedData, CertifiedGaps.ReseedData.Valid,
    CertifiedGaps.ReseedData.spatial, CertifiedGaps.thermalPurity_ge_of_dominant}
  \leanok
  \uses{lem:seed-capped, lem:seed-temperature-lift}
  Let \(R\) be valid reseeding data, \(m=ak\), and suppose \(Z\) admits arbitrarily accurate
  finite real power sums at \(\{k,m,2m,n,2n\}\).
  If \(1<Z_m\), \((1-p)Z_m^2\le Z_{2m}\), \(Z_k\le B\), and \(Z_n>0\), then
  \((1-\es)Z_n^2\le Z_{2n}\), with the parameters of \(R\).
  For a finite energy family at length \(N\), assume \(\beta>0\), some \(E_i<-\kappa N\),
  \(\widetilde Z_N(\beta)\le q\), \(R_\tau\ge0\), \((q-1)^2<R_\tau\),
  \(p_\tau\le1\), and \((1-p_\tau)(1+R_\tau)^2\le1\).
  Then \(\Pt(N,2\beta)\ge1-p_\tau\).
\end{theorem}
\begin{proof}
  \leanok
  The old purity forces a dominant power-sum entry; the short trace controls its relative tail.
  Pass to limits. Apply the dominant Boltzmann-weight bound for the thermal conclusion.
\end{proof}

\section{Transfer matrices and polynomial filters}

Odd short-ring lengths require a single-site Trotter construction.
Its finite symmetric matrices replace the continuous-history operator of
\cite[Section 3]{OpenAIHaldane}; only the resulting trace inequalities pass to the limit.

\begin{lemma}[Single-site transfer matrices at all lengths]
  \label{lem:single-site-transfer}
  \lean{CertifiedGaps.SingleSite.transfer, CertifiedGaps.SingleSite.transfer_isHermitian,
    CertifiedGaps.SingleSite.trace_stepOp_pow,
    CertifiedGaps.SingleSite.tendsto_stepOp_pow,
    CertifiedGaps.SingleSite.eventually_abs_twisted_transfer_sub_le}
  \leanok
  \uses{lem:chernoff}
  For real bond factors \(C_\alpha\) and real on-site matrix \(o\), use the symmetrized site
  factor \(M(a,b)=(C_aC_b+C_bC_a)/2\), with \(C_{\varnothing}=I\).
  The time-label transfer matrix \(X_M\) is real symmetric.
  For nonnegative time step, the trace of the \(M\)-fold step on a ring is
  \(\sum_\gamma D_r(\gamma)(X_M^L)_{\pi\gamma,\gamma}\), for closing-bond signs \(r\) and
  relabeling \(\pi\).
  At step \(\beta/M\), the physical product converges to \(e^{-\beta H_L(r,\pi)}\) for
  \(L\ge2\). For \(\beta\ge0\), the transfer traces therefore approximate all twists on any
  prescribed finite set of such lengths.
\end{lemma}
\begin{proof}
  \leanok
  Exchange the space and time products. The step has value \(I\) and derivative \(-H_L\) at zero,
  so the finite-dimensional product limit applies.
\end{proof}

\begin{lemma}[Sign and cyclic sectors]
  \label{lem:single-site-sectors}
  \lean{CertifiedGaps.SingleSite.SignSymmetry, CertifiedGaps.SingleSite.LabelSymmetry,
    CertifiedGaps.SingleSite.D2Symmetry.d2Approx,
    CertifiedGaps.SingleSite.A4Symmetry.sectorApprox_ion,
    CertifiedGaps.exists_cyclic_sectors}
  \leanok
  \uses{lem:single-site-transfer, def:seed-scheme}
  Suppose two trace-preserving sign symmetries and a cyclic relabeling satisfy
  \texttt{D2Symmetry}. At \(\beta\ge0\), the plain and sign-twisted traces have the \(D_2\)
  sector approximation on every finite set of lengths at least \(2\).
  If the relabeling has order dividing three, as required by \texttt{A4Symmetry}, the additional
  cyclic twist gives the three-sector approximation.
  The two nonreal characters of the cyclic action have equal spectra by complex conjugation.
\end{lemma}
\begin{proof}
  \leanok
  Sign characters make the transfer matrix block diagonal. The cycle identifies the three
  nontrivial \(D_2\) blocks and splits its trivial block into the cyclic sectors.
\end{proof}

\begin{lemma}[Polynomial caps from moments]
  \label{lem:seed-filter}
  \lean{CertifiedGaps.PolynomialFilter.CapCheck,
    CertifiedGaps.PolynomialFilter.pairingBound,
    CertifiedGaps.PolynomialFilter.abs_lt_of_filter_bounds}
  \leanok
  Let \(F\) be a rational polynomial nonnegative on \(\R\), and let a finite family \(x_i\)
  satisfy the stated lower and upper moment bounds.
  Suppose \texttt{pairingBound} returns \(B\), and \texttt{CapCheck} verifies that
  \(F(x)>B\) on both rays \(|x|\ge v\).
  Then every member of the family satisfies \(|x_i|<v\).
\end{lemma}
\begin{proof}
  \leanok
  Nonnegativity gives \(F(x_i)\le\sum_jF(x_j)\le B\).
  The rational Taylor-coefficient check excludes both rays.
\end{proof}

\begin{lemma}[Periodic matrix-product trial bounds]
  \label{lem:mps-variational}
  \lean{CertifiedGaps.MPS.state, CertifiedGaps.MPS.sum_state_sq,
    CertifiedGaps.MPS.dotProduct_ringHamiltonian_mulVec,
    CertifiedGaps.MPS.exp_lt_gibbsTrace_int,
    CertifiedGaps.exists_eigenvalue_lt_of_re_dotProduct_lt}
  \leanok
  \uses{def:ring-operators}
  For integer tensors \(A_s\), put
  \(\psi(\sigma)=\Tr(A_{\sigma_0}\cdots A_{\sigma_{N-1}})\),
  \(T=\sum_sA_s\otimes A_s\), and
  \(O=\sum h_{st,s't'}A_sA_t\otimes A_{s'}A_{t'}\).
  For a symmetric integer bond \(h\) and \(N\ge2\),
  \(\|\psi\|^2=\Tr T^N\) and
  \(\langle\psi,H_N\psi\rangle=N\Tr(T^{N-2}O)\).
  Thus, for integers \(a>0\) and \(b\) with \(\kappa=b/a\), the inequality
  \(a\Tr(T^{N-2}O)<-b\Tr T^N\) implies an energy below \(-\kappa N\)
  and \(Z_N(\beta)>e^{\beta\kappa N}\) for \(\beta>0\).
\end{lemma}
\begin{proof}
  \leanok
  Contract the doubled virtual indices and use cyclicity for each bond.
  Apply the variational principle and retain the corresponding Boltzmann weight.
\end{proof}

\begin{lemma}[Open matrix-product trial bounds]
  \label{lem:open-mps-variational}
  \lean{CertifiedGaps.OpenMPS.state, CertifiedGaps.OpenMPS.sum_state_sq,
    CertifiedGaps.OpenMPS.dotProduct_openChainHamiltonian_mulVec,
    CertifiedGaps.OpenMPS.exists_eigenvalue_lt}
  \leanok
  \uses{lem:mps-variational}
  Use integer tensors and end vectors to form
  \(\psi(\sigma)=b^{\mathsf T}A_{\sigma_0}\cdots A_{\sigma_{N-1}}c\).
  With \(L=b\otimes b\) and \(u=c\otimes c\), its norm is \(LT^Nu\).
  Each bond contributes \(LT^jOT^{N-2-j}u\); end terms contribute
  \(LG_0T^{N-1}u\) and \(LT^{N-1}G_1u\), where
  \(G_i=\sum_{s,t}(f_i)_{st}A_s\otimes A_t\).
  For symmetric integer bonds and end matrices, with real field coefficients, a contraction
  inequality \(\langle\psi,H\psi\rangle<E\|\psi\|^2\) gives an eigenvalue below \(E\).
\end{lemma}
\begin{proof}
  \leanok
  Keep the virtual endpoints open in the contraction identities and apply the variational
  principle.
\end{proof}

\section{Open-chain initialization and the leading eigenvalue}

The boundary argument needs the leading spectral line and its overlap with a boundary vector.
The finite-matrix construction retains these data through the Trotter limit.

\begin{lemma}[Open traces in the single-site construction]
  \label{lem:single-site-open}
  \lean{CertifiedGaps.SingleSite.capVec, CertifiedGaps.SingleSite.trace_openStep_pow,
    CertifiedGaps.SingleSite.tendsto_openPairing}
  \leanok
  \uses{lem:single-site-transfer}
  Give the left and right ends coefficients \(u_L,u_R\), representing
  \(-\sum_\alpha u_{L,\alpha}C_\alpha\) and
  \(-\sum_\alpha u_{R,\alpha}C_\alpha\).
  For a chain of \(k+1\) sites and nonnegative time step, its Trotter trace is the bilinear
  pairing \(v_L^{\mathsf T}X_M^{k+1}v_R\), where each cap vector is the product of its
  time-label weights. At step \(\beta/M\), the real parts of these pairings converge to the open
  Gibbs trace
  for \(\beta\ge0\).
\end{lemma}
\begin{proof}
  \leanok
  Contract the two end labels into the cap vectors and use the derivative of the open step
  at zero.
\end{proof}

\begin{definition}[A leading line and its symmetry projection]
  \label{def:leading-line}
  \lean{CertifiedGaps.IsLeadingLine, CertifiedGaps.SingleSite.avgProj,
    CertifiedGaps.SingleSite.avgProj_conjTranspose,
    CertifiedGaps.SingleSite.avgProj_mul_self,
    CertifiedGaps.SingleSite.isLeadingLine_of_blocks}
  \leanok
  \uses{lem:single-site-sectors, def:spectral-forms}
  A leading-line witness consists of a Hermitian matrix \(X\), a unit vector \(\psi\),
  \(X\psi=e_0\psi\), and
  \(|q_{X-e_0|\psi\rangle\langle\psi|}(y)|\le r\|y\|^2\).
  For the \(A_4\) action, the average projection is
  \(\Pi=\tfrac14(I+D_x)(I+D_y)\tfrac13(I+C+C^2)\).
  It is an orthogonal projection. Scale the transfer matrix by \(t>0\). Suppose one eigenvalue
  of the trivial \(D_2\) block has modulus \(|e_0|>r\ge0\), while the other eigenvalues of that
  block and all eigenvalues of one nontrivial block \(N\) have modulus at most \(r\).
  Then the scaled transfer matrix has a leading line with eigenvalue \(e_0\); its vector lies in
  the trivial \(D_2\) block and is fixed by the cyclic relabeling.
\end{definition}

\begin{lemma}[Boundary moments and a leading-weight polynomial]
  \label{lem:leading-polynomial}
  \lean{CertifiedGaps.IsLeadingLine.openMomentWitness,
    CertifiedGaps.IsLeadingLine.mul_poly_le}
  \leanok
  \uses{def:leading-line, def:open-moment-bounds}
  Suppose a leading-line witness has \(r\le|e_0|\), and a unitary \(U\) commutes with \(X\)
  and fixes \(\psi\).
  The even ring moments \(\Tr X^{2N}\), even pairings
  \(\operatorname{Re}\langle v,X^{2N}v\rangle\), and odd pairings
  \(\operatorname{Re}\langle v,X^{2N+1}Uv\rangle\), for \(N\ge1\), have an open-moment
  witness with \(\mu=|e_0|\), \(c=|\langle\psi,v\rangle|^2\),
  \(s=\operatorname{sign}(e_0)\) (taking \(s=1\) at zero), and \(\rho\mu\le r\).
  If an orthogonal projection \(P\) commutes with \(X\) and fixes \(\psi\), and
  \(f(x)=\sum_{n\in S}a_nx^n\ge0\) on \(\R\), with every \(n\in S\) positive, then
  \[
    c f(e_0)\le\sum_{n\in S}a_n\operatorname{Re}\langle v,X^nPv\rangle.
  \]
\end{lemma}
\begin{proof}
  \leanok
  Split off the leading rank-one projection. The remaining spectral weights are nonnegative
  in the even moments and in the projected polynomial.
\end{proof}

\begin{lemma}[Leading witnesses pass to limits]
  \label{lem:leading-limit}
  \lean{CertifiedGaps.LeadingWitness, CertifiedGaps.LeadingWitness.of_tendsto,
    CertifiedGaps.LeadingWitness.mul_poly_le}
  \leanok
  \uses{def:open-moment-bounds}
  A \texttt{LeadingWitness} adds \(\ell\le\mu\), \(\rho\le\rho_{\max}\), and
  \(c\sum_{n\in S}a_n(s\mu)^n\le\sum_{n\in S}a_nI_n\) to an open-moment witness.
  Suppose such witnesses exist eventually along approximants, whose ring and open traces
  converge at every length at least \(2\), and whose \(I_n\) converge for \(n\in S\).
  If the limiting \(Z_2,Z_4\) are positive, a limiting witness exists.
  If also \(\ell>0\), \(\rho_{\max}<1\), all odd \(G_{2N+1}>0\) for \(N\ge1\), and
  \(\sum_{n\in S}a_nI_n\le B\) with \(B\ge0\), then \(c\sum_{n\in S}a_n\mu^n\le B\).
\end{lemma}
\begin{proof}
  \leanok
  The witnesses lie in a compact box controlled by \(Z_2,Z_4,G_2\).
  Pass the closed inequalities along a convergent subsequence.
  Positivity of odd traces forces \(c=0\) or \(s=1\).
  This finite-matrix limit replaces the compact-operator argument of
  \cite[Sections 2--3 and Appendix A]{OpenAIBoundary}.
\end{proof}

\begin{lemma}[A decidable ratio bound]
  \label{lem:ratio-check}
  \lean{CertifiedGaps.PolynomialFilter.RatioCheck,
    CertifiedGaps.PolynomialFilter.pow_div_sq_le_of_ratioCheck}
  \leanok
  Suppose \(\texttt{RatioCheck}(P,k_0,v)\) holds, \(v\ge0\), \(1\le k\), and \(k_0\le k\).
  For \(v\le x\le U\),
  \[
    \frac{x^{2k}}{P(x)^2}\le\frac{U^{2k}}{P(U)^2}.
  \]
\end{lemma}
\begin{proof}
  \leanok
  The rational coefficient checks imply \(P(x)>0\) and \(xP'(x)\le k_0P(x)\) on the ray.
  Differentiate the ratio.
\end{proof}

\begin{lemma}[Initializing a boundary defect and thermal purity]
  \label{lem:open-initialization}
  \lean{CertifiedGaps.OpenMomentWitness.residual_le_of_bounds,
    CertifiedGaps.OpenMomentWitness.one_sub_mul_le_of_ground,
    CertifiedGaps.OpenMomentWitness.le_add_of_bounds,
    CertifiedGaps.OpenMomentWitness.boundaryDefect_le,
    CertifiedGaps.thermalPurity_ge_of_ground}
  \leanok
  \uses{def:open-moment-bounds, lem:ratio-check}
  Let an open-moment witness have \(0<\mu\le U\), \(\rho\le\rho'\),
  \(G_{2K}\le\Gamma\), and \(1\le K\le N\). Then
  \[
    G_{2N}-c\mu^{2N}\le(\rho'U)^{2N-2K}\Gamma.
  \]
  If this upper bound is at most \(\epsilon g\), with \(0<g\le G_{2N}\), then
  \((1-\epsilon)G_{2N}\le c\mu^{2N}\).
  When \(Z_4>0\), \(\rho<1\), and \(G_k>0\), any bound
  \((1-E)G_k\le c\mu^k\) implies boundary defect at most \(E\).
  For a finite energy family with \(\beta_0\ge0\), some \(E_i\le e_0\), \(R\ge0\), and
  \(Z(\beta_0)\le e^{-\beta_0e_0}(1+R)\), its thermal purity at \(2^J\beta_0\)
  is at least \((1+R^{2^J})^{-2}\).
  With the same witness bounds, if also \(cP(\mu)^2\le C\),
  \(\texttt{RatioCheck}(P,k_0,v)\) holds, \(0\le v\le\mu\), and \(k_0\le N\), then
  \(G_{2N}\le CU^{2N}/P(U)^2+(\rho'U)^{2N-2K}\Gamma\).
\end{lemma}
\begin{proof}
  \leanok
  Use the even residual decay and the trial-state lower bound.
  A bound on \(cP(\mu)^2\), together with Lemma~\ref{lem:ratio-check}, also bounds
  the leading part of \(G_{2N}\). Compare powers of Boltzmann weights for the last claim.
\end{proof}
